% Iconos de los diagramas, dibujados en TikZ.
%
% No se usa ningun paquete de iconos: fontawesome y compania anadian una
% dependencia que descargar y una tipografia ajena a la plantilla. Se dibujan
% a mano, con los mismos colores que el resto de las figuras, y asi cumplen
% ademas la regla del tutor de que las imagenes sencillas se hacen y no se
% copian.
%
% Cada icono es un `tikzpicture` que se puede meter DENTRO de un nodo de otro
% `tikzpicture`. El parametro opcional es la escala; por defecto 0.22. Por
% debajo de 4 mm en la pagina el detalle se pierde al reducir la figura al
% ancho de la caja de texto, asi que todos son siluetas y no dibujos.
%
% TODOS declaran la MISMA caja envolvente con `\useasboundingbox`. Sin ella
% cada icono ocupa lo que ocupa su trazo y los nodos que los contienen salen
% de alturas distintas.
%
% Solo estan los que se usan. Si hace falta uno nuevo, se anade aqui y queda
% disponible para todas las figuras a la vez.

\tikzset{
  iconoTrazo/.style={draw=diagNodoBorde, line width=0.9pt,
    line join=round, line cap=round},
}

% Caja comun. Se llama al principio de cada icono.
\newcommand{\iconocaja}{\useasboundingbox (-1.4,-1.35) rectangle (1.4,1.35);}

% --- La web de la que se parte ------------------------------------------
\newcommand{\icPantalla}[1][0.22]{%
  \begin{tikzpicture}[scale=#1, iconoTrazo, every path/.style={iconoTrazo}]
    \iconocaja
    \draw[fill=white, rounded corners=0.14] (-1.25,-0.95) rectangle (1.25,1.0);
    \draw (-1.25,0.45) -- (1.25,0.45);
    \foreach \d in {-0.98,-0.68,-0.38}{\fill[diagNodoBorde] (\d,0.72) circle (0.09);}
    \draw (-0.9,0.05) -- (0.9,0.05);
    \draw (-0.9,-0.35) -- (0.55,-0.35);
    \draw (-0.9,-0.7) -- (0.25,-0.7);
  \end{tikzpicture}}

% --- La extraccion: la web, y de ella se baja el contenido --------------
\newcommand{\icWeb}[1][0.22]{%
  \begin{tikzpicture}[scale=#1, iconoTrazo, every path/.style={iconoTrazo}]
    \iconocaja
    \draw[fill=white] (0,0.4) circle (0.8);
    \draw (-0.8,0.4) -- (0.8,0.4);
    \draw (0,1.2) arc(90:270:0.4 and 0.8);
    \draw (0,1.2) arc(90:-90:0.4 and 0.8);
    \draw[-{Latex[length=1.6mm]}, line width=1.2pt] (0,-0.65) -- (0,-1.25);
  \end{tikzpicture}}

% --- Cortar en fragmentos -----------------------------------------------
% Ensena el RESULTADO del corte y no la herramienta: unas tijeras a 5 mm se
% leen como un conector. Tres losas desplazadas se reconocen aunque el icono
% sea diminuto.
\newcommand{\icCortar}[1][0.22]{%
  \begin{tikzpicture}[scale=#1, iconoTrazo, every path/.style={iconoTrazo}]
    \iconocaja
    \draw[fill=white] (-0.95,0.5)  rectangle (0.75,1.2);
    \draw[fill=white] (-0.7,-0.35) rectangle (1.0,0.35);
    \draw[fill=white] (-0.95,-1.2) rectangle (0.75,-0.5);
  \end{tikzpicture}}

% --- La base de datos vectorial -----------------------------------------
\newcommand{\icBase}[1][0.22]{%
  \begin{tikzpicture}[scale=#1, iconoTrazo, every path/.style={iconoTrazo}]
    \iconocaja
    \draw[fill=white] (-0.85,0.85) -- (-0.85,-0.85) arc(180:360:0.85 and 0.32)
    -- (0.85,0.85) arc(0:180:0.85 and 0.32) -- cycle;
    \draw (-0.85,0.28) arc(180:360:0.85 and 0.32);
    \draw (-0.85,-0.3) arc(180:360:0.85 and 0.32);
  \end{tikzpicture}}

% --- Buscar / recuperar -------------------------------------------------
\newcommand{\icLupa}[1][0.22]{%
  \begin{tikzpicture}[scale=#1, iconoTrazo, every path/.style={iconoTrazo}]
    \iconocaja
    \draw[fill=white] (0.2,0.4) circle (0.78);
    \draw[line width=1.5pt] (-0.35,-0.15) -- (-1.05,-1.15);
  \end{tikzpicture}}

% --- El modelo de lenguaje ----------------------------------------------
\newcommand{\icChip}[1][0.22]{%
  \begin{tikzpicture}[scale=#1, iconoTrazo, every path/.style={iconoTrazo}]
    \iconocaja
    \draw[fill=white] (-0.8,-0.8) rectangle (0.8,0.8);
    \draw (-0.42,-0.42) rectangle (0.42,0.42);
    \foreach \d in {-0.45,0,0.45}{
      \draw (-0.8,\d) -- (-1.2,\d);
      \draw (0.8,\d) -- (1.2,\d);
      \draw (\d,0.8) -- (\d,1.2);
      \draw (\d,-0.8) -- (\d,-1.2);}
  \end{tikzpicture}}

% --- La respuesta, en forma de conversacion -----------------------------
\newcommand{\icGlobo}[1][0.22]{%
  \begin{tikzpicture}[scale=#1, iconoTrazo, every path/.style={iconoTrazo}]
    \iconocaja
    \draw[fill=white, rounded corners=0.28] (-1.1,-0.3) rectangle (1.1,1.2);
    % Se borra el trozo de borde por donde sale el rabo y se redibujan sus dos
    % lados: dibujar el rabo antes dejaba la linea del globo cruzandolo.
    \path[fill=white, draw=none] (-0.62,-0.27) -- (-0.16,-0.27) -- (-0.5,-1.15) -- cycle;
    \draw (-0.62,-0.3) -- (-0.5,-1.15) -- (-0.16,-0.3);
    \draw (-0.6,0.45) -- (0.6,0.45);
    \draw (-0.6,0.8) -- (0.35,0.8);
  \end{tikzpicture}}

% --- Quien pregunta -----------------------------------------------------
\newcommand{\icPersona}[1][0.22]{%
  \begin{tikzpicture}[scale=#1, iconoTrazo, every path/.style={iconoTrazo}]
    \iconocaja
    \draw[fill=white] (0,0.85) circle (0.45);
    \draw (0,0.4) -- (0,-0.5);
    \draw (-0.7,0.05) -- (0.7,0.05);
    \draw (0,-0.5) -- (-0.55,-1.25);
    \draw (0,-0.5) -- (0.55,-1.25);
  \end{tikzpicture}}

% --- Una decision escrita (ADR) -----------------------------------------
\newcommand{\icDocumento}[1][0.22]{%
  \begin{tikzpicture}[scale=#1, iconoTrazo, every path/.style={iconoTrazo}]
    \iconocaja
    \draw[fill=white] (-0.75,-1.2) -- (-0.75,1.2) -- (0.3,1.2) -- (0.8,0.7)
    -- (0.8,-1.2) -- cycle;
    \draw (0.3,1.2) -- (0.3,0.7) -- (0.8,0.7);
    \draw (-0.45,0.25) -- (0.5,0.25);
    \draw (-0.45,-0.2) -- (0.5,-0.2);
    \draw (-0.45,-0.65) -- (0.2,-0.65);
  \end{tikzpicture}}

% --- Un instrumento de medida (verificador o banco) ---------------------
\newcommand{\icLista}[1][0.22]{%
  \begin{tikzpicture}[scale=#1, iconoTrazo, every path/.style={iconoTrazo}]
    \iconocaja
    \draw[fill=white] (-0.8,-1.2) rectangle (0.8,1.05);
    \draw[fill=white] (-0.35,0.85) rectangle (0.35,1.28);
    \draw (-0.5,0.3) -- (-0.3,0.1) -- (0.0,0.5);
    \draw (0.2,0.3) -- (0.55,0.3);
    \draw (-0.5,-0.4) -- (-0.3,-0.6) -- (0.0,-0.2);
    \draw (0.2,-0.4) -- (0.55,-0.4);
  \end{tikzpicture}}
